% ============================================================================
% Globalisation and the Welfare State: Evidence from 32 OECD Economies,
% 1980–2023
%
% Scaffolded manuscript. Sections below are placeholders to be filled in
% as results converge. Regression tables and figures generated by the
% pipeline are imported via \input from outputs/tables/ — do not edit
% those files by hand; re-run `make data` or the relevant exporter in
% src/analysis/robustness.py and src/analysis/latex_injector.py.
% ============================================================================

\documentclass[11pt,a4paper]{article}

% --- Packages -----------------------------------------------------------------
\usepackage[utf8]{inputenc}
\usepackage[T1]{fontenc}
\usepackage[british]{babel}
\usepackage{amsmath,amssymb}
\usepackage{booktabs}
\usepackage{graphicx}
\usepackage{siunitx}
\usepackage[margin=1in]{geometry}
\usepackage[hidelinks]{hyperref}
\usepackage{natbib}

\graphicspath{{../outputs/figures/}}

% --- Metadata -----------------------------------------------------------------
\title{Globalisation and the Welfare State: \\
  Evidence from 32 OECD Economies, 1980--2023}
\author{Anton Ebsen\thanks{%
  Replication materials:
  \url{https://github.com/AntonEbsen/economics-of-the-welfare-state}.}}
\date{\today}

\begin{document}
\maketitle

\begin{abstract}
  \noindent This paper investigates whether deeper integration into the global
  economy erodes the redistributive capacity of the welfare state. Using
  a balanced panel of 32 OECD countries over 1980--2023, I estimate
  two-way fixed-effects regressions of social security transfers on the
  KOF Globalisation Index and its economic, social, and political
  sub-components. I document a \emph{(sign to be filled in)} association
  and probe its robustness across welfare-regime subgroups, subperiod
  splits, and an event-study design around country-level structural
  breaks.
\end{abstract}

% ============================================================================
\section{Introduction}
\label{sec:intro}
% ============================================================================

% TODO: motivation, literature gap, contribution, roadmap.

% ============================================================================
\section{Related Literature}
\label{sec:literature}
% ============================================================================

% TODO: Compensation hypothesis (Rodrik, Katzenstein), efficiency
% hypothesis (Swank, Garrett), recent panel evidence.

% ============================================================================
\section{Data}
\label{sec:data}
% ============================================================================

The dataset merges six public sources into a balanced country-year
panel of 32 OECD economies over 1980--2023 (1{,}408 observations). The
dependent variable, \emph{social security transfers as a percentage of
GDP} (\texttt{sstran}), is drawn from the Comparative Political Data
Set \citep{armingeon2024cpds}. Globalisation is measured by the KOF
Globalisation Index \citep{gygli2019kof}, with separate entries for the
economic, social, and political sub-indices. Standard controls come
from World Bank WDI: real GDP per capita, consumer price inflation,
total population, and the age dependency ratio.

A complete description of sources, variable construction, and missing
data patterns appears in \texttt{paper/methods.md}. Verification of raw
input checksums is enforced by \texttt{scripts/download\_raw\_data.py}.

% ============================================================================
\section{Empirical Strategy}
\label{sec:strategy}
% ============================================================================

The baseline specification is a two-way fixed-effects regression
\begin{equation}
  \mathrm{sstran}_{it} = \alpha_i + \lambda_t
    + \beta \cdot \mathrm{KOFGI}_{i,t-1}
    + \mathbf{X}_{i,t-1}^{\top} \boldsymbol{\gamma} + \varepsilon_{it},
  \label{eq:baseline}
\end{equation}
where $\alpha_i$ and $\lambda_t$ are country and year fixed effects,
$\mathbf{X}_{i,t-1}$ collects the lagged controls enumerated above, and
standard errors are clustered at the country level.

I supplement \eqref{eq:baseline} with
(i) subperiod splits (pre- versus post-2008) to test for regime
change,
(ii) heterogeneity by welfare-regime type
\citep{espingandersen1990three}, and
(iii) an event-study decomposition around country-specific structural
breaks detected via the Chow/QLR procedures in
\texttt{src/clean/structural\_breaks.py}.

% ============================================================================
\section{Results}
\label{sec:results}
% ============================================================================

\subsection{Baseline}
% TODO: narrative; table below is auto-imported.
% \input{../outputs/tables/baseline_regression.tex}

\subsection{Stepwise robustness}
% TODO: narrative.
% \input{../outputs/tables/stepwise_robustness.tex}

\subsection{Subperiod and regime heterogeneity}
% TODO: narrative.
% \input{../outputs/tables/subperiod_regressions.tex}

\subsection{Event study}
% TODO: reference figure.
% \begin{figure}[t]
%   \centering
%   \includegraphics[width=0.8\linewidth]{event_study.pdf}
%   \caption{Event study around structural breaks.}
% \end{figure}

% ============================================================================
\section{Discussion}
\label{sec:discussion}
% ============================================================================

% TODO: interpretation, policy implications, limitations.

% ============================================================================
\section{Conclusion}
\label{sec:conclusion}
% ============================================================================

% TODO: headline finding, contribution, avenues for future work.

% ============================================================================
\bibliographystyle{apalike}
\bibliography{bibliography}

\appendix
% ============================================================================
% Appendix: additional tables, diagnostics, and specification tests.
% Included from manuscript.tex via % ============================================================================
% Appendix: additional tables, diagnostics, and specification tests.
% Included from manuscript.tex via % ============================================================================
% Appendix: additional tables, diagnostics, and specification tests.
% Included from manuscript.tex via \input{appendix}.
% ============================================================================

\section{Data Appendix}
\label{app:data}

\subsection{Country coverage}
The 32 countries in the balanced panel are listed in
\texttt{paper/methods.md}. Five of them (BGR, EST, LTU, LVA, SVK) enter
the dataset from 1995 owing to WB data availability; the merge fills
forward these rows with NaN for earlier years, and the two-way
fixed-effects estimator drops them via within-transformation.

\subsection{Missing-data patterns}
% TODO: Import the missing-data heatmap emitted by src/clean/quality.py
% when the narrative is finalised.

% ============================================================================
\section{Additional Results}
\label{app:results}
% ============================================================================

\subsection{First-stage diagnostics}
% TODO: Hausman test results, serial correlation tests (Wooldridge),
% cross-sectional dependence (Pesaran CD). These are emitted by
% src/analysis/regression_utils.py::run_hausman_test and exported via
% src/clean/stats.py.

\subsection{Alternative specifications}
% TODO: Random-effects, first-differenced, dynamic panel (Arellano-Bond)
% comparison.

\subsection{Placebo distribution}
% TODO: Histogram of placebo β from run_placebo_test (shuffles the
% treatment within entity). Exported to outputs/figures/placebo_hist.pdf.

\subsection{Structural-break detection}
% TODO: Chow and QLR statistics per country from
% src/clean/structural_breaks.py.

% ============================================================================
\section{Replication Notes}
\label{app:replication}
% ============================================================================

The full replication archive --- raw inputs, cleaning pipeline,
regression code, and this manuscript --- is available at
\url{https://github.com/AntonEbsen/economics-of-the-welfare-state}.
See \texttt{docs/replication\_notes.md} for system requirements,
runtime expectations, and reproduction instructions.
.
% ============================================================================

\section{Data Appendix}
\label{app:data}

\subsection{Country coverage}
The 32 countries in the balanced panel are listed in
\texttt{paper/methods.md}. Five of them (BGR, EST, LTU, LVA, SVK) enter
the dataset from 1995 owing to WB data availability; the merge fills
forward these rows with NaN for earlier years, and the two-way
fixed-effects estimator drops them via within-transformation.

\subsection{Missing-data patterns}
% TODO: Import the missing-data heatmap emitted by src/clean/quality.py
% when the narrative is finalised.

% ============================================================================
\section{Additional Results}
\label{app:results}
% ============================================================================

\subsection{First-stage diagnostics}
% TODO: Hausman test results, serial correlation tests (Wooldridge),
% cross-sectional dependence (Pesaran CD). These are emitted by
% src/analysis/regression_utils.py::run_hausman_test and exported via
% src/clean/stats.py.

\subsection{Alternative specifications}
% TODO: Random-effects, first-differenced, dynamic panel (Arellano-Bond)
% comparison.

\subsection{Placebo distribution}
% TODO: Histogram of placebo β from run_placebo_test (shuffles the
% treatment within entity). Exported to outputs/figures/placebo_hist.pdf.

\subsection{Structural-break detection}
% TODO: Chow and QLR statistics per country from
% src/clean/structural_breaks.py.

% ============================================================================
\section{Replication Notes}
\label{app:replication}
% ============================================================================

The full replication archive --- raw inputs, cleaning pipeline,
regression code, and this manuscript --- is available at
\url{https://github.com/AntonEbsen/economics-of-the-welfare-state}.
See \texttt{docs/replication\_notes.md} for system requirements,
runtime expectations, and reproduction instructions.
.
% ============================================================================

\section{Data Appendix}
\label{app:data}

\subsection{Country coverage}
The 32 countries in the balanced panel are listed in
\texttt{paper/methods.md}. Five of them (BGR, EST, LTU, LVA, SVK) enter
the dataset from 1995 owing to WB data availability; the merge fills
forward these rows with NaN for earlier years, and the two-way
fixed-effects estimator drops them via within-transformation.

\subsection{Missing-data patterns}
% TODO: Import the missing-data heatmap emitted by src/clean/quality.py
% when the narrative is finalised.

% ============================================================================
\section{Additional Results}
\label{app:results}
% ============================================================================

\subsection{First-stage diagnostics}
% TODO: Hausman test results, serial correlation tests (Wooldridge),
% cross-sectional dependence (Pesaran CD). These are emitted by
% src/analysis/regression_utils.py::run_hausman_test and exported via
% src/clean/stats.py.

\subsection{Alternative specifications}
% TODO: Random-effects, first-differenced, dynamic panel (Arellano-Bond)
% comparison.

\subsection{Placebo distribution}
% TODO: Histogram of placebo β from run_placebo_test (shuffles the
% treatment within entity). Exported to outputs/figures/placebo_hist.pdf.

\subsection{Structural-break detection}
% TODO: Chow and QLR statistics per country from
% src/clean/structural_breaks.py.

% ============================================================================
\section{Replication Notes}
\label{app:replication}
% ============================================================================

The full replication archive --- raw inputs, cleaning pipeline,
regression code, and this manuscript --- is available at
\url{https://github.com/AntonEbsen/economics-of-the-welfare-state}.
See \texttt{docs/replication\_notes.md} for system requirements,
runtime expectations, and reproduction instructions.


\end{document}
